\documentclass[12pt,a4paper]{amsart}

\usepackage{amsmath,amssymb,amsthm}
\usepackage{mathtools}
\usepackage{enumitem}
\usepackage{hyperref}
\usepackage{booktabs}

%%%─ Environments
\newtheorem{theorem}{Theorem}[section]
\newtheorem{lemma}[theorem]{Lemma}
\newtheorem{proposition}[theorem]{Proposition}
\newtheorem{corollary}[theorem]{Corollary}
\newtheorem{conjecture}[theorem]{Conjecture}
\theoremstyle{definition}
\newtheorem{definition}[theorem]{Definition}
\newtheorem{assumption}[theorem]{Assumption}
\newtheorem{problem}[theorem]{Open Problem}
\theoremstyle{remark}
\newtheorem{remark}[theorem]{Remark}

%%%─ Macros
\newcommand{\Kc}{\mathcal{K}_c}
\newcommand{\Tr}{\mathrm{Tr}}
\newcommand{\lam}[2]{\lambda_{#2}^{(#1)}}
\newcommand{\cA}{\mathcal{A}}
\newcommand{\Ai}{\mathrm{Ai}}
\newcommand{\Bi}{\mathrm{Bi}}
\newcommand{\R}{\mathbb{R}}
\newcommand{\C}{\mathbb{C}}
\newcommand{\op}{\mathrm{op}}
\newcommand{\tr}{\mathrm{tr}}

\subjclass[2020]{47B34, 45C05, 35P20, 47A55, 34E20, 60G55}
\keywords{Prolate spheroidal wave functions, sinc kernel, resolvent asymptotics,
Airy operator, double-scaling limit, spectral counting, Fredholm determinant}

\begin{document}

\title[Paper XIV: Airy Resolvent for Sinc Kernel]{%
 Paper~XIV: Airy Resolvent Asymptotics for the Sinc Kernel Operator\\[0.3em]
 \large An Honest Gap Analysis}

\author{\textit{prolate-gram-coercivity programme}}
\date{May 2026}

\begin{abstract}
This paper provides a rigorous account of what is and what is \emph{not}
established in the programme to prove condition~(Q5) via resolvent
asymptotics.

\medskip\noindent
\textbf{What this paper does.}
\begin{enumerate}[nosep]
 \item Derives the double-scaling structure $(c, c^{2/3}, c^{1/3})$
 from the PSWE turning-point analysis (unconditional).
 \item Identifies the \emph{correct} target: not the deterministic
 count $N_\varphi(c)=2$, but the spectral measure
 convergence statement~(Q5$''$) below.
 \item Introduces the Bridge Lemma (Lemma~\ref{lem:bridge})
 as a conditional implication: $L^2$ kernel convergence with rate
 $\Rightarrow$ trace-norm resolvent convergence.
 \item Provides a rigorous regularisation framework for $\det(\cA-\zeta)$.
 \item Explicitly lists five open sub-problems whose resolution
 would complete the programme.
\end{enumerate}

\medskip\noindent
\textbf{What this paper does NOT do.}
\begin{itemize}[nosep]
 \item It does \emph{not} prove that the sinc resolvent converges
 to the Airy resolvent in trace norm (open: Problem~\ref{prob:bridge}).
 \item It does \emph{not} prove that the bulk term $T_0$ is holomorphic
 in $\zeta$ (open: Problem~\ref{prob:bulksep}).
 \item It does \emph{not} prove $N_\varphi(c)=2$ (this requires all of
 the above plus a matching argument, Problem~\ref{prob:matching}).
\end{itemize}

\medskip\noindent
\textbf{Revised target (Q5$''$).}
For a smooth test function $\varphi \in C_c^\infty(\R)$,
\begin{equation*}
 \Tr[\varphi(c^{1/3}(\Kc - z_{n^*}))] \;\longrightarrow\;
 \Tr[\varphi(\cA)] \quad \text{as } c\to\infty,
\end{equation*}
where $\cA = -d^2/d\xi^2 + \xi$ on $L^2(\R_+)$.
This is spectral measure convergence in the double-scaling variable.
The binary count $N_\varphi(c)=2$ follows as a \emph{secondary consequence}
only after an additional matching argument (Section~\ref{sec:matching}).
\end{abstract}

\maketitle
\tableofcontents
\bigskip\hrule\bigskip

%% ============================================================
\section{The Three Scales: Unconditional Derivation}
\label{sec:scales}
%% ============================================================

The prolate spheroidal wave equation (PSWE) is
\begin{equation}\label{eq:PSWE}
 \frac{d}{dx}\bigl[(1-x^2)\psi'\bigr] + [\chi_n - c^2 x^2]\psi = 0,
 \quad x\in[-1,1],
\end{equation}
with turning point $x_0 = \sqrt{\chi_n}/c$.
In the transition regime $n^* \sim 2c/\pi$,
the turning point approaches $x_0 \to 1^-$.

\begin{lemma}[Scale derivation --- unconditional]
\label{lem:scales}
Let $x = 1 - c^{-2/3}\xi$ with $\xi = O(1)$, and write
$\chi_n = c^2 - 2c^{4/3}\tau + O(c^{2/3})$ where
$\tau = (n - 2c/\pi)c^{-1/3}/\pi^{-1}$.
Then equation~\eqref{eq:PSWE} transforms to
\begin{equation}\label{eq:Airyform}
 \frac{d^2\psi}{d\xi^2} = (\xi - \tau)\psi + O(c^{-1/3})
\end{equation}
in leading order as $c\to\infty$.
Consequently:
\begin{enumerate}[nosep,label=(\roman*)]
 \item the spatial scale of the transition region is $c^{-2/3}$;
 \item the spectral scale (eigenvalue spacing near $n^*$) is $c^{-1/3}$;
 \item the bulk density scale is $c$ (Weyl: $\Tr[\Kc] \sim 2c/\pi$).
\end{enumerate}
These are the unique scales compatible with~\eqref{eq:Airyform};
no other scales appear at leading order.
\end{lemma}

\begin{proof}
Standard computation: insert $x = 1-c^{-2/3}\xi$ into~\eqref{eq:PSWE},
expand $(1-x^2) = 2c^{-2/3}\xi + O(c^{-4/3})$
and $c^2 x^2 = c^2 - 2c^{4/3}\xi + O(c^{2/3})$.
Dividing by the leading coefficient gives~\eqref{eq:Airyform}.
\end{proof}

\begin{remark}[This is the only unconditional result of Paper~XIV]
Lemma~\ref{lem:scales} is proved without further assumptions.
All remaining results are conditional (conjectural or placeholder).
\end{remark}

%% ============================================================
\section{The Correct Target: Spectral Measure Convergence}
\label{sec:target}
%% ============================================================

\begin{definition}[Scaled spectral measure]
\label{def:specmeas}
For $c > 0$, define the scaled spectral measure of $\Kc$ near $n^*$:
\begin{equation}\label{eq:specmeas}
 \mu_c := \sum_j \delta_{c^{1/3}(\lam{c}{j} - z_{n^*})},
\end{equation}
where $z_{n^*} = \lam{c}{n^*}$ is the transition eigenvalue.
\end{definition}

\begin{remark}[Why $N_\varphi(c)=2$ is unstable as a primary target]
\label{rem:fragility}
The binary claim $N_\varphi(c) = 2$ requires:
\begin{enumerate}[nosep,label=(F\arabic*)]
 \item \emph{Centering:} $z_{n^*}$ lies between the first two Airy
 eigenvalues $\{-a_k\}$ in the $\zeta$-variable.
 \item \emph{Window calibration:} $\zeta_0$ is chosen so that
 $[-\zeta_0,\zeta_0]$ contains $\{-a_1,-a_2\}$ but not $-a_3$.
 For the Airy zeros ($a_1\approx 2.34$, $a_2\approx 4.09$,
 $a_3\approx 5.52$), this requires $\zeta_0 \in (2.04,\,2.76)$;
 it is a finite-width interval, not an asymptotic regime.
 \item \emph{Matching:} the transition index $n^*$ corresponds to the
 correct Airy eigenvalue label $k$.
 This is a sub-leading Weyl-law condition on $n^* - 2c/\pi$
 that is \emph{not} determined by leading-order asymptotics alone.
\end{enumerate}
None of (F1)--(F3) is automatic.
The claim $N_\varphi(c)=2$ is therefore a consequence,
not a starting point, of the programme.
\end{remark}

\begin{conjecture}[Spectral measure convergence --- primary target]
\label{conj:specmeas}
For every $\varphi \in C_c^\infty(\R)$,
\begin{equation}\label{eq:Q5pp}
 \int \varphi\,d\mu_c
 \;=\; \Tr[\varphi(c^{1/3}(\Kc - z_{n^*}))]
 \;\longrightarrow\;
 \Tr[\varphi(\cA)]
 \;=\; \sum_k \varphi(-a_k)
 \quad (c\to\infty).
\end{equation}
This is condition~\textbf{(Q5$''$)} of this paper.
\end{conjecture}

\begin{remark}[Relation to Airy point process]
Conjecture~\ref{conj:specmeas} is a \emph{vague topology} convergence
of measures, weaker than weak convergence of point processes.
It is the minimal condition needed to extract counting information.
The Airy point process convergence (Paper~XIV-v1 Conj.\ 1.3)
is a strictly stronger statement and is not assumed here.
\end{remark}

\subsection{From spectral measure to binary count: the matching problem}
\label{sec:matching}

\begin{problem}[Matching and binary count]
\label{prob:matching}
Assume Conjecture~\ref{conj:specmeas}.
To conclude $N_\varphi(c) = 2$ for large $c$, one additionally needs:
\begin{enumerate}[nosep,label=(M\arabic*)]
 \item \emph{Uniform rate:}
 $|\int\varphi\,d\mu_c - \sum_k\varphi(-a_k)| = O(c^{-\alpha})$
 for some $\alpha>0$, with constant independent of $\varphi$.
 \item \emph{Index matching:}
 The index $n^*$ satisfies $n^* - 2c/\pi = c^{1/3}\tau_0 + o(c^{1/3})$
 for a specific $\tau_0$ determined by the first Airy zero $a_1$.
 This requires sub-leading Weyl asymptotics for $\Kc$.
 \item \emph{Window calibration (F2) above.}
\end{enumerate}
(M2) is a \emph{new open problem} not present in Papers~I--XIII.
\end{problem}

%% ============================================================
\section{The Resolvent Programme}
\label{sec:resolvent}
%% ============================================================

\subsection{The Helffer--Sj\"ostrand approach}

For $\varphi \in C_c^\infty(\R)$ with almost-analytic extension $\tilde\varphi$:
\begin{equation}\label{eq:HS}
 \Tr[\varphi(c^{1/3}(\Kc - z_{n^*}))]
 = -\frac{1}{\pi}\int_\C \partial_{\bar\zeta}\tilde\varphi(\zeta)\,
 \Tr[(\Kc - z_{n^*} - c^{-1/3}\zeta)^{-1}]\,d^2\zeta.
\end{equation}
For this to converge to $\Tr[\varphi(\cA)]$, one needs:
\begin{equation}\label{eq:traceconv}
 \Tr[(\Kc - z_{n^*} - c^{-1/3}\zeta)^{-1}]
 \;\longrightarrow\;
 c^{1/3}\Tr[(\cA - \zeta)^{-1}]
 \quad \text{as } c\to\infty,
\end{equation}
locally uniformly for $\zeta \in \C\setminus\R$,
at least in a suitable integrated-in-$\zeta$ sense.

\begin{remark}[The bulk--edge separation problem]
\label{rem:bulkedge}
Let $z = z_{n^*} + c^{-1/3}\zeta$.
The full resolvent trace decomposes heuristically as
\begin{equation}\label{eq:decomp}
 \Tr[(\Kc - z)^{-1}]
 \sim c\cdot T_0(z) + c^{2/3}\cdot \Tr[(\cA-\zeta)^{-1}] + \ldots
\end{equation}
The term $c\cdot T_0(z)$ represents the bulk contribution
(eigenvalues far from $z_{n^*}$).
\textbf{Critical gap:}
$T_0(z)$ is \emph{not automatically holomorphic in $\zeta$}.
Writing $z = z_{n^*} + c^{-1/3}\zeta$,
the function $z \mapsto T_0(z)$ is meromorphic in $z$
(poles at all bulk eigenvalues),
but its dependence on $\zeta$ enters both through $z$ explicitly
and through the implicit dependence of $z_{n^*}$ on $c$.
The claim that the $c\cdot T_0$ term vanishes from~\eqref{eq:HS}
after integrating against $\partial_{\bar\zeta}\tilde\varphi$
requires a \emph{microlocal separation argument}:
the bulk resolvent must be shown to be analytic (in $\zeta$)
after localisation to the edge region.
This is the content of Problem~\ref{prob:bulksep}.
\end{remark}

\begin{problem}[Bulk--edge microlocal separation]
\label{prob:bulksep}
Define a microlocal cutoff $\chi_c$ separating bulk
($|\lam{c}{j} - z_{n^*}| \gg c^{-1/3}$) from edge contributions.
Prove:
\begin{enumerate}[nosep,label=(B\arabic*)]
 \item The bulk part $T_0^{(\mathrm{cut})}(\zeta)
 := \sum_{j:\,|\lam{c}{j}-z_{n^*}|>Rc^{-1/3}}
 (\lam{c}{j} - z_{n^*} - c^{-1/3}\zeta)^{-1}$
 is holomorphic in $\zeta$ for $|\zeta| < R/2$.
 \item Its contribution to~\eqref{eq:HS} is $O(c^{-1/3})$
 after integration against $\partial_{\bar\zeta}\tilde\varphi$.
\end{enumerate}
\end{problem}

\subsection{The Bridge Lemma}

\begin{lemma}[Bridge Lemma --- conditional]
\label{lem:bridge}
Let $B_c := c^{1/3}(\Kc - z_{n^*})|_{\mathrm{edge}}$ be the
edge-localised scaled operator (formal; defined up to Problem~\ref{prob:bulksep}).
Suppose:
\begin{enumerate}[label=(BR\arabic*),nosep]
 \item The operator $B_c$ converges in \emph{strong resolvent sense}
 to $\cA$ as $c\to\infty$.
 \item The family $(B_c - \zeta)^{-1}$ is uniformly trace-norm
 bounded for $\zeta$ in compact subsets of $\C\setminus\R$,
 with bound $\|(B_c-\zeta)^{-1}\|_{\tr} \leq C(K)|\mathrm{Im}\,\zeta|^{-1}$.
 \item The difference $(B_c - \cA)(\cA - \zeta)^{-2}$ is trace class
 with $\|(B_c - \cA)(\cA-\zeta)^{-2}\|_{\tr} = O(c^{-\alpha})$
 for some $\alpha > 0$.
\end{enumerate}
Then:
\begin{equation}\label{eq:bridgeconclusion}
 \bigl|\Tr[(B_c - \zeta)^{-1}] - \Tr[(\cA-\zeta)^{-1}]\bigr|
 = O(c^{-\alpha}|\mathrm{Im}\,\zeta|^{-2})
\end{equation}
and Conjecture~\ref{conj:specmeas} holds for $\varphi \in C_c^\infty$
with support in a compact subset of $\R$.
\end{lemma}

\begin{proof}[Proof (conditional on (BR1)--(BR3))]
The resolvent identity gives
$(B_c-\zeta)^{-1} - (\cA-\zeta)^{-1}
 = (\cA-\zeta)^{-1}(\cA-B_c)(B_c-\zeta)^{-1}$.
Taking trace norms and applying (BR3):
$\|(B_c-\zeta)^{-1}-(\cA-\zeta)^{-1}\|_{\tr}
\leq \|(\cA-B_c)(\cA-\zeta)^{-2}\|_{\tr}\cdot
\|(B_c-\zeta)^{-1}\|_{\op}\cdot\|(\cA-\zeta)\|_{\op}$,
which gives~\eqref{eq:bridgeconclusion}.
Plugging into~\eqref{eq:HS} and using $|\partial_{\bar\zeta}\tilde\varphi| \leq C_N|\mathrm{Im}\,\zeta|^N$
gives the trace convergence.
\end{proof}

\begin{remark}[Status of conditions (BR1)--(BR3)]
Condition (BR1) may follow from form convergence of the PSWE
to the Airy equation (Lemma~\ref{lem:scales}), using Kato's
theory of convergence of quadratic forms~\cite{Kato1966}.
Conditions (BR2) and (BR3) are the genuinely new requirements:
(BR2) requires a uniform Schatten-class bound for the edge resolvent,
and (BR3) requires that the \emph{difference} of the two operators
(sinc edge minus Airy) is trace class after regularisation.
Neither is standard and neither follows from (BR1) alone.
This is the analytic core of Paper~XIV.
\end{remark}

%% ============================================================
\section{The Fredholm Determinant Route}
\label{sec:fredholm}
%% ============================================================

\subsection{Regularisation of $\det(\cA - \zeta)$}

\begin{remark}[The Airy determinant requires regularisation]
\label{rem:detreg}
Since $\cA$ is unbounded with $\sigma(\cA) = \{-a_k\} \to +\infty$,
the expression $\det(\cA-\zeta)$ is not a Fredholm determinant
in the classical sense.
We use the \emph{relative Fredholm determinant}:
\begin{equation}\label{eq:reldet}
 D_\cA^{\mathrm{rel}}(\zeta;\zeta_0)
 := \det\bigl[(\cA-\zeta)(\cA-\zeta_0)^{-1}\bigr],
\end{equation}
where $\zeta_0 \notin \sigma(\cA)$ is a fixed reference point.
Since $(\cA-\zeta)(\cA-\zeta_0)^{-1} = I + (\zeta_0-\zeta)(\cA-\zeta_0)^{-1}$
and $(\cA-\zeta_0)^{-1}$ is trace class
(as $\cA$ has discrete spectrum with $\sum_k a_k^{-p} < \infty$ for $p>3/2$),
$D_\cA^{\mathrm{rel}}(\zeta;\zeta_0)$ is a genuine Fredholm determinant
in the sense of Simon~\cite{Simon2005}.
Its zeros in $\zeta$ are exactly the eigenvalues $\{-a_k\}$,
independent of $\zeta_0$.
\end{remark}

\subsection{Fredholm conjecture and Rouch\'{e} argument}

\begin{conjecture}[Fredholm convergence --- conditional on bulk separation]
\label{conj:fredholm}
Assume Problem~\ref{prob:bulksep} is resolved.
Let $D_{\mathrm{edge}}^{(c)}(\zeta)
:= \det[(B_c - \zeta)(B_c - \zeta_0)^{-1}]$
(relative determinant of the edge operator).
Then
\begin{equation}\label{eq:fredconv}
 D_{\mathrm{edge}}^{(c)}(\zeta)
 \;\longrightarrow\;
 D_\cA^{\mathrm{rel}}(\zeta;\zeta_0)
\end{equation}
uniformly on compact subsets of $\C$,
with rate $O(c^{-\alpha})$ (same $\alpha$ as in Bridge Lemma).
\end{conjecture}

\begin{remark}[Rouch\'{e} argument --- conditional]
If Conjecture~\ref{conj:fredholm} holds,
then by operator-valued Rouch\'{e} (Gohberg--Krein~\cite{GohKre1969}),
for any contour $\Gamma$ enclosing a finite set of Airy eigenvalues
and avoiding all others,
the number of zeros of $D_{\mathrm{edge}}^{(c)}$ inside $\Gamma$
equals the number of zeros of $D_\cA^{\mathrm{rel}}$ inside $\Gamma$
for $c$ large enough.
This gives a finite count of PSWF eigenvalues near $z_{n^*}$
\emph{in terms of Airy zeros},
but still requires the index matching (Problem~\ref{prob:matching})
to conclude that exactly two eigenvalues lie in a given $c^{-1/3}$-window.
\end{remark}

%% ============================================================
\section{Open Problems: Honest Summary}
\label{sec:open}
%% ============================================================

\begin{problem}[Strong resolvent convergence (BR1)]
\label{prob:BR1}
Prove that the edge-scaled operator $B_c$ converges to $\cA$
in strong resolvent sense.
Approach: show quadratic form convergence
$\mathfrak{q}[B_c](u,u) \to \mathfrak{q}[\cA](u,u)$
for $u \in H^1_0(\R_+)$, using Lemma~\ref{lem:scales}.
\end{problem}

\begin{problem}[Trace-class difference (BR3) --- core problem]
\label{prob:bridge}
Prove $\|(B_c - \cA)(\cA-\zeta)^{-2}\|_{\tr} = O(c^{-\alpha})$.
This requires:
\begin{enumerate}[nosep,label=(\roman*)]
 \item $L^2$-kernel convergence of $\Kc^{(\mathrm{sc})}$ to $P_{(-\infty,0]}(\cA)$
 with rate;
 \item upgrading $L^2$ kernel convergence to Schatten-$p$ convergence;
 \item applying Simon's interpolation~\cite{Simon2005} to obtain trace-class.
\end{enumerate}
\emph{This is the deepest analytic problem in Paper~XIV.}
\end{problem}

\begin{problem}[Bulk--edge separation]
\label{prob:bulksep2}
See Problem~\ref{prob:bulksep} above.
Remarks: for the sinc kernel on $[-1,1]$,
bulk and edge are not microlocally separated by support
(unlike for the sine kernel on $\R$).
A Fourier restriction argument or pseudodifferential
calculus adapted to $[-1,1]$ is needed.
\end{problem}

\begin{problem}[Sub-leading Weyl law and index matching]
\label{prob:matching2}
See Problem~\ref{prob:matching}.
Specifically: prove that
\[
 n^* - \frac{2c}{\pi} = c^{1/3}\tau_0 + o(c^{1/3})
\]
where $\tau_0$ is determined by the Airy zero $a_1$ and the phase
condition of the PSWE at the turning point.
This is a sub-leading asymptotic for the PSWF count;
not known from current literature.
\end{problem}

\begin{problem}[(Q5-lower) via quasimode injection]
\label{prob:lower}
Under Conjecture~\ref{conj:specmeas},
prove the lower bound $N_\varphi(c) \geq 2 - O(c^{-\alpha})$
by constructing two approximate eigenvectors
$f_{n^*}^{(c)}, f_{n^*+1}^{(c)}$ in the Airy regime
with residual $O(c^{-1/3-\varepsilon})$.
This is the \emph{accessible} direction of (Q5).
\end{problem}

%% ============================================================
\section{Programme State After Paper~XIV}
\label{sec:state}
%% ============================================================

\begin{center}
\renewcommand{\arraystretch}{1.35}
\begin{tabular}{p{4.5cm} p{1.8cm} p{1.5cm}}
\toprule
\textbf{Claim} & \textbf{Type} & \textbf{Status}\\
\midrule
Three scales from PSWE (Lem.~\ref{lem:scales}) & deductive & \checkmark\\
Relative Airy det.\ well-defined (Rem.~\ref{rem:detreg}) & algebraic & \checkmark\\
Bulk holomorphicity not automatic (Rem.~\ref{rem:bulkedge}) & structural & \checkmark\\
Bridge Lemma form (Lem.~\ref{lem:bridge}) & conditional & \checkmark\\
$N=2$ fragility identified (Rem.~\ref{rem:fragility}) & structural & \checkmark\\
Strong resolvent conv.\ (BR1) (Prob.~\ref{prob:BR1}) & analytic & open\\
Trace-class difference (BR3) (Prob.~\ref{prob:bridge}) & analytic & open (hardest)\\
Bulk--edge separation (Prob.~\ref{prob:bulksep2}) & analytic & open\\
Index matching (Prob.~\ref{prob:matching2}) & Weyl theory & open (new!)\\
(Q5-lower) quasimode (Prob.~\ref{prob:lower}) & constructive & open (accessible)\\
\bottomrule
\end{tabular}
\end{center}

\medskip\noindent
\textbf{Logical chain for Gap-S:}
\[
 \underbrace{\text{(BR1)+(BR3)}+\text{bulk sep.}+\text{matching}}_{
 \text{four open problems}}
 \;\Longrightarrow\;
 \text{(Q5$''$)}
 \;\Longrightarrow\;
 N_\varphi(c)=2
 \;\Longrightarrow\;
 \text{Gap-S \emph{(Thm.~B, Paper~XIII).}}
\]

\begin{thebibliography}{99}
\bibitem{paper13} \textit{Paper~XIII: Obstructions to Gap-S},
 \texttt{paper13\_gap\_s.tex}, 2026.
\bibitem{Kato1966} T.~Kato, \textit{Perturbation Theory for Linear Operators},
 Springer, 1966.
\bibitem{Osipov2013} A.~Osipov, V.~Rokhlin, H.~Xiao,
 \textit{Prolate Spheroidal Wave Functions of Order Zero}, Springer, 2013.
\bibitem{Widom1964} H.~Widom,
 \textit{Asymptotic behavior of the eigenvalues of certain integral equations},
 Trans.\ Amer.\ Math.\ Soc.\ \textbf{109} (1964), 278--295.
\bibitem{TracyWidom1994} C.~Tracy, H.~Widom,
 \textit{Level-spacing distributions and the Airy kernel},
 Comm.\ Math.\ Phys.\ \textbf{159} (1994), 151--174.
\bibitem{DKMVZ1999} P.~Deift, T.~Kriecherbauer, K.~McLaughlin,
 S.~Venakides, X.~Zhou,
 \textit{Uniform asymptotics for polynomials orthogonal with respect
 to varying exponential weights},
 Comm.\ Pure Appl.\ Math.\ \textbf{52} (1999), 1335--1425.
\bibitem{Simon2005} B.~Simon,
 \textit{Trace Ideals and Their Applications}, 2nd ed., AMS, 2005.
\bibitem{GohKre1969} I.~Gohberg, M.~Krein,
 \textit{Introduction to the Theory of Linear Nonselfadjoint Operators},
 AMS, 1969.
\bibitem{HelSjo1989} B.~Helffer, J.~Sj\"ostrand,
 \textit{\'Equation de Schr\"odinger avec champ magn\'etique},
 Lecture Notes in Physics \textbf{345}, Springer, 1989.
\end{thebibliography}

\end{document}
